% !TEX root = ../algebra_02.tex

\section{Применение теоремы о гомоморфизме к классификации циклических групп}

\begin{Definition}{Циклическая группа}{def:cyclic}
	Группа $G$ называется \emph{циклической}, если существует элемент $g\in G$ такой, что
	\[
	G=\langle g\rangle = \{g^k\mid k\in \ZZ\}.
	\]
	Элемент $g$ называется \emph{образующим} группы $G$.
\end{Definition}

\begin{Theorem}{Все подгруппы группы $\ZZ$}{subgroups-Z}
	Всякая подгруппа $H\le \ZZ$ имеет вид
	\[
	H=n\ZZ=\langle n\rangle
	\]
	для некоторого $n\ge 0$.
\end{Theorem}

\begin{proof}
	Если $H=\{0\}$, то утверждение верно при $n=0$.

	Пусть теперь $H\ne\{0\}$. Тогда в $H$ есть положительные элементы. Пусть $n$ --- наименьший положительный элемент из $H$. Покажем, что $H=n\ZZ$. Во-первых, $n\ZZ\subseteq H$, поскольку $H$ замкнута относительно сложения и взятия противоположного.

	Во-вторых, пусть $x\in H$. По алгоритму деления существует представление
	\[
	x=qn+r, \qquad 0\le r<n.
	\]
	Так как $x\in H$ и $qn\in H$, получаем $r=x-qn\in H.$
	По минимальности $n$ имеем $r=0$. Значит, $x\in n\ZZ$.

	Следовательно, $H=n\ZZ$.
\end{proof}

\begin{Theorem}{Классификация циклических групп}{thm:cyclic-classification}
	Пусть $G=\langle g\rangle$ --- циклическая группа.
	\begin{enumerate}
		\item Если $|G|=\infty$, то $G\cong \ZZ$.
		\item Если $|G|=n<\infty$, то $G\cong \ZZ/n\ZZ$.
	\end{enumerate}
\end{Theorem}

\begin{proof}
	Рассмотрим гомоморфизм
	\[
	\varphi\colon \ZZ \to G, \qquad k\mapsto g^k.
	\]
	Он сюръективен, поскольку $G=\langle g\rangle$. По теореме о подгруппах группы $\ZZ$ существует $n\ge 0$ такое, что $\Ker\varphi=n\ZZ.$

	Если $n=0$, то $\Ker\varphi=\{0\}$, и потому $\varphi$ инъективен. Следовательно, $G\cong \ZZ$. Если $n>0$, то по первой теореме об изоморфизме
	\[
	\ZZ/n\ZZ = \ZZ/\Ker\varphi \cong \Im\varphi = G.
	\]
	Если при этом $|G|<\infty$, то обязательно $n=|G|$.
\end{proof}
